\section{Experiments and Analysis}
\label{sec:experiments}

\subsection{Comparison with Frontier Models}
\label{sec:exp:baselines}

To contextualize our fine-tuned 8B model, we evaluate three frontier models on a 500-question sample from v4 and on the 54-question independent holdout (Table~\ref{tab:baselines}).

\begin{table}[t]
\centering
\small
\caption{Frontier and fine-tuned model accuracy on reliability engineering benchmarks.}
\label{tab:baselines}
\begin{tabular}{@{}lcc@{}}
\toprule
Model & v4 (500q) & Holdout (54q) \\
\midrule
Claude Sonnet 4.6      & \textbf{92.4\%} & \textbf{81.1\%} \\
o3-mini                & 86.4\% & 77.1\% \\
Gemini 3.1 Pro         & 76.4\% & 82.6\%$^\dagger$ \\
\midrule
Qwen3-8B (base)        & 63.6\%$^*$ & 53.7\% \\
Qwen3-8B (SFT)         & 82.2\%$^*$ & 33.3\% \\
Qwen3-8B (SFT+GRPO)   & 57.9\%$^{**}$ & 53.7\% \\
\bottomrule
\end{tabular}
\vspace{0.1em}

{\scriptsize $^*$10-fold CV. $^{**}$On 266 mixed-difficulty questions. $^\dagger$23/54 extracted.}
\end{table}

Our SFT model (82.2\% under 10-fold CV) places between o3-mini and Gemini~3.1~Pro despite being an order of magnitude smaller. On the holdout, SFT regresses sharply to 33.3\% ($-20.4$pp, $p=0.019$), while GRPO fully recovers to 53.7\%, matching the base model (Section~\ref{sec:grpo:results}).

\subsection{Analysis}
\label{sec:exp:analysis}

\paragraph{Dataset size scaling.} SFT improvement scales non-linearly with dataset size (Table~\ref{tab:sft_results}): marginal below 300 examples, a threshold effect around 500, and strong gains at 866 ($+18.6\%$, $p=0.002$). The question flip ratio grows from 1.3:1 (215 samples) to 4.6:1 (866 samples), indicating that larger augmented datasets produce more robust improvements with fewer catastrophic regressions.

\paragraph{Epoch convergence.} Early stopping (patience 2) consistently selects epoch 4 across all dataset sizes and configurations, confirming a stable convergence point.

\paragraph{GRPO overfitting.} GRPO accuracy degrades with continued training past step~80 ($-3.8$pp at step~100, $-5.6$pp at step~200), suggesting 266 questions are insufficient for sustained improvement. The truncation rate increases with training (4 for base, 10 for GRPO step~80), indicating GRPO encourages longer reasoning chains.
